\documentclass[12pt,a4paper]{article}
\usepackage[T2A]{fontenc}
\usepackage[utf8]{inputenc}
\usepackage[russian]{babel}
\usepackage{float}
\usepackage{amsmath, amssymb}
\usepackage{geometry}
\usepackage{tikz}
\usepackage{tikz-3dplot}
\usepackage{caption}
\usetikzlibrary{arrows.meta, calc, decorations.pathmorphing}
\geometry{margin=2cm}

\newcommand{\vu}{\bar u}
\newcommand{\vv}{\bar v}
\newcommand{\vy}{\bar y}
\newcommand{\va}{\bar a}
\newcommand{\vb}{\bar b}
\newcommand{\vc}{\bar c}
\newcommand{\vd}{\bar d}
\newcommand{\ve}{\bar e}
\newcommand{\vo}{\bar o}

\begin{document}

\setcounter{page}{107}

\begin{center}
  \textbf{ЛЕКЦИЯ 15. ЛИНЕЙНАЯ ЗАВИСИМОСТЬ И ЛИНЕЙНАЯ НЕЗАВИСИМОСТЬ ВЕКТОРОВ}
\end{center}

\subsection*{\textbf{15.1. Понятие линейной зависимости и линейной
независимости}}

Пусть дано линейное пространство $V$ и в нем система векторов
$\{\vu_1, \ldots, \vu_m\}$.

\textbf{Определение.} Линейной комбинацией векторов $\vu_1, \vu_2,
\ldots, \vu_m \in V$ называется вектор вида
\begin{equation}
  \vy = \lambda_1 \vu_1 + \lambda_2 \vu_2 + \ldots + \lambda_m \vu_m, \tag{15.1}
\end{equation}
где $\lambda_1, \lambda_2, \ldots, \lambda_m$ --- произвольные числа.

\textbf{Определение.} Линейная комбинация (15.1) называется
тривиальной, если $\lambda_1 = \lambda_2 = \ldots = \lambda_m = 0$. В
противном случае линейная комбинация называется нетривиальной (т.е.
если хотя бы одно $\lambda_i \ne 0$).

\textbf{Определение.} Система векторов $\Sigma = \{\vu_1, \vu_2,
\ldots, \vu_m\}$ называется линейно зависимой, если существует
нетривиальная линейная комбинация этих векторов, равная нулевому
вектору, т.е. если существуют числа $\lambda_1, \lambda_2, \ldots,
\lambda_m$, одновременно не равные нулю, такие что:
$$
\lambda_1 \vu_1 + \lambda_2 \vu_2 + \ldots + \lambda_m \vu_m = \vo.
$$

\textbf{Определение.} Система векторов называется линейно
независимой, если нулевому вектору равна только тривиальная линейная
комбинация этих векторов.

Замечание. Краткое определение:

Векторы $\vu_1, \vu_2, \ldots, \vu_m$ --- линейно независимы, если
$$
\lambda_1 \vu_1 + \lambda_2 \vu_2 + \ldots + \lambda_m \vu_m = \vo
\Leftrightarrow \lambda_1 = \lambda_2 = \ldots = \lambda_m = 0.
$$
В противном случае векторы линейно зависимы.

\subsection*{\textbf{15.2. Свойства линейно зависимых и линейно
независимых систем векторов}}

\textbf{1.} Система из одного вектора линейно зависима тогда и только
тогда, когда этот вектор нулевой.

\textbf{Доказательство.} Линейная зависимость системы из одного
вектора равносильна тому, что $\lambda \vu = \vo$, а условие $\lambda
\vu = \vo \Leftrightarrow \vu = \vo$ ($\lambda \ne 0$) --- следствие
простейших свойств линейных пространств.

\textbf{ч.т.д.}

\textbf{2.} Критерий линейной зависимости системы векторов.

Система векторов $\Sigma = \{\vu_1, \vu_2, \ldots, \vu_m\}$, где $m >
1$, линейно зависима тогда и только тогда, когда хотя бы один из
векторов этой системы линейно выражается через другие.

\textbf{Доказательство}

1) Необходимость

Если $\Sigma = \{\vu_1, \vu_2, \ldots, \vu_m\}$ линейно зависима $\Rightarrow$

$\exists$ числа $\lambda_1, \lambda_2, \ldots, \lambda_m$,
одновременно не равные нулю и такие, что
\begin{equation}
\lambda_1 \vu_1 + \lambda_2 \vu_2 + \ldots + \lambda_m \vu_m = \vo \tag{15.2}
\end{equation}

Пусть $\lambda_i \ne 0$. Тогда в силу (15.2)
$$
\vu_i = -\frac{\lambda_1}{\lambda_i} \vu_1 -
\frac{\lambda_2}{\lambda_i} \vu_2 - \cdots -
\frac{\lambda_{i-1}}{\lambda_i} \vu_{i-1} -
\frac{\lambda_{i+1}}{\lambda_i} \vu_{i+1} - \cdots -
\frac{\lambda_m}{\lambda_i} \vu_m
$$

2) Достаточность

Пусть, например, $\vu_1 = \lambda_2 \vu_2 + \ldots + \lambda_m \vu_m$. Тогда
$$
\vo = -\vu_1 + \lambda_2 \vu_2 + \ldots + \lambda_m \vu_m
$$
--- это нетривиальная линейная комбинация $\Rightarrow \Sigma =
\{\vu_1, \vu_2, \ldots, \vu_m\}$ линейно зависима.

\textbf{ч.т.д.}

\textbf{3.} Если подсистема системы векторов линейно зависима, то и
вся система линейно зависима.

\textbf{4.} Любая подсистема линейно независимой системы векторов
линейно независима.

\textbf{5.} Критерий линейной независимости системы векторов. Система
векторов $\Sigma = \{\vu_1, \vu_2, \ldots, \vu_m\}$ линейно
независима тогда и только тогда, когда любой вектор, являющийся
линейной комбинацией этих векторов, имеет единственное разложение по
этим векторам.

Доказательство элементарное. Доказывается методом от противного.

\textbf{6.} Если система векторов $\Sigma = \{\vu_1, \vu_2, \ldots,
\vu_m\}$ --- линейно независима, а система векторов $\Sigma' =
\{\vu_1, \vu_2, \ldots, \vu_m, \vv\}$ линейно зависима, то вектор
$\vv$ линейно выражается через векторы $\vu_1, \vu_2, \ldots, \vu_m$.

\textbf{Доказательство}

$\Sigma'$ --- линейно зависима $\Rightarrow$
\begin{equation}
\lambda_1 \vu_1 + \lambda_2 \vu_2 + \ldots + \lambda_m \vu_m +
\lambda_0 \vv = \vo, \tag{15.3}
\end{equation}
где хотя бы одно из чисел $\lambda_1, \lambda_2, \ldots, \lambda_m,
\lambda_0$ не равно нулю.

Если $\lambda_0 = 0$, то ненулевой коэффициент $\lambda_k$ находится
среди чисел $\lambda_1, \lambda_2, \ldots, \lambda_m$.

При этом равенство (15.3) переходит в равенство
\begin{equation}
\lambda_1 \vu_1 + \lambda_2 \vu_2 + \ldots + \lambda_m \vu_m = \vo, \tag{15.4}
\end{equation}

где (15.4) --- нетривиальная линейная комбинация, равная нулевому
вектору $\Rightarrow$ противоречие линейной независимости системы
$\Sigma = \{\vu_1, \vu_2, \ldots, \vu_m\}$ $\Rightarrow \lambda_0 \ne
0 \Rightarrow$ вектор $\vv$ линейно выражается из равенства (15.3)
через векторы $\vu_1, \vu_2, \ldots, \vu_m$.

\textbf{ч.т.д.}

\textbf{Пример}

Рассмотрим в $\mathbb{R}^n$ ---мерном пространстве арифметических
векторов векторы:
\begin{align*}
\ve_1 &= (1, 0, \ldots, 0);\\
\ve_2 &= (0, 1, \ldots, 0);\\
&\ldots\\
\ve_n &= (0, 0, \ldots, 1).
\end{align*}

$\Sigma = \{\ve_1, \ve_2, \ldots, \ve_n\}$ --- линейно независима,
т.к. линейная комбинация этих векторов с коэффициентами $\lambda_1,
\lambda_2, \ldots, \lambda_n$ представляет собой арифметический
вектор $(\lambda_1, \lambda_2, \ldots, \lambda_n)$, который равен
$\vo = (0, \ldots, 0) \Leftrightarrow \lambda_i = 0$, $i = 1, 2, \ldots, n$.

\textbf{Пример}

Пространство многочленов

Многочлены $1, t, t^2, \ldots, t^n$ линейно независимы, т.к. их
линейная комбинация с коэффициентами $\lambda_0, \lambda_1,
\lambda_2, \ldots, \lambda_n$ представляет собой многочлен, который
равен нулю тогда и только тогда, когда все $\lambda_i = 0$, $i = 0,
1, \ldots, n$.

\subsection*{\textbf{15.3. Геометрический смысл линейной зависимости
и линейной независимости для системы геометрических векторов}}

\textbf{Теорема.} Два вектора линейно зависимы тогда и только тогда,
когда они коллинеарны.

\textbf{Доказательство}

1) необходимость;

Если векторы $\va$ и $\vb$ линейно зависимы $\Rightarrow$ согласно
критерию линейной зависимости
$$
\va = \lambda \vb \Rightarrow \va \parallel \vb.
$$

2) достаточность.

Пусть $\va \parallel \vb$ и $\va \ne \vo$ (если $\va = \vo$, то
согласно свойствам \textbf{1)} и \textbf{3)} $\va$ и $\vb$ линейно зависимы)
$\Rightarrow \vb = \lambda \va \Rightarrow$ согласно свойству
\textbf{2)} $\va$ и $\vb$ линейно зависимы.

\textbf{ч.т.д.}

\textbf{Теорема.} Три вектора линейно зависимы тогда и только тогда,
когда они компланарны.

\textbf{Доказательство}

1) необходимость;

Пусть векторы $\va, \vb, \vc$ линейно зависимы $\Rightarrow$ один из
них линейно выражается через другие. Пусть $\vc = \alpha \va + \beta \vb$.

Если $\va \parallel \vb \Rightarrow$ векторы $\va, \vb, \vc$
коллинеарны и тем более компланарны.

Если $\va \nparallel \vb$. Тогда приведем векторы $\va$ и $\vb$ к
одному началу. Вектор $\vc$ является диагональю параллелограмма,
построенного на векторах $\alpha \va$ и $\beta \vb$ $\Rightarrow \vc$
лежит в той же плоскости, что и векторы $\va$ и $\vb$ $\Rightarrow$
векторы $\va, \vb, \vc$ --- компланарны

2) достаточность.

Пусть $\va, \vb, \vc$ --- компланарны и $\va \nparallel \vb$ (если
$\va \parallel \vb \Rightarrow \va$ и $\vb$ --- линейно зависимы
$\Rightarrow$ $\va, \vb, \vc$ --- линейно зависимы согласно свойству \textbf{3}.

Приведем векторы $\va, \vb, \vc$ к одному началу

\begin{figure}[H]
\centering
\begin{tikzpicture}[scale=1.5]
\coordinate (O) at (0,0);
\coordinate (A) at (1.5,1);
\coordinate (B) at (2.5,0);
\coordinate (C) at (4,1);

\draw[->,thick] (O) -- (A) node[midway,left] {$\va$};
\draw[->,thick] (O) -- (B) node[midway,below] {$\vb$};
\draw[->,thick] (O) -- (C) node[midway,above left] {$\vc$};
\draw[dashed] (A) -- (C);
\draw[dashed] (B) -- (C);

\node[below left] at (O) {$O$};
\node[above] at (A) {$A$};
\node[below] at (B) {$B$};
\node[above right] at (C) {$C$};
\end{tikzpicture}
\caption*{Рисунок 15.1}
\end{figure}

$\vc = \overline{OA} + \overline{OB} = \alpha \va + \beta \vb
\Rightarrow \vc$ линейно выражается через $\va$ и $\vb$ $\Rightarrow
\va, \vb, \vc$ линейно зависимы.

\textbf{ч.т.д.}

\textbf{Теорема.} Любые четыре вектора линейно зависимы.

\textbf{Доказательство}

$\Sigma = \{\va, \vb, \vc, \vd\}$ --- линейно зависима.

Рассмотрим $\Sigma' = \{\va, \vb, \vc\}$

1) Если $\Sigma' = \{\va, \vb, \vc\}$ линейно зависима $\Rightarrow
\Sigma = \{\va, \vb, \vc, \vd\}$ линейно зависима.

2) Если $\Sigma' = \{\va, \vb, \vc\}$ линейно независима $\Rightarrow
\va, \vb, \vc$ --- некомпланарны.

\begin{figure}[H]
\centering
\begin{tikzpicture}[scale=1.5]
\coordinate (O) at (0,0);
\coordinate (A) at (2,0);
\coordinate (B) at (0.5,1.5);
\coordinate (C) at (1,2.5);
\coordinate (D) at (3,2.5);

\draw[->,thick] (O) -- (A) node[midway,below] {$\va$};
\draw[->,thick] (O) -- (B) node[midway,left] {$\vb$};
\draw[->,thick] (O) -- (C) node[midway,left] {$\vc$};
\draw[->,thick] (O) -- (D) node[midway,above left] {$\vd$};

\draw[dashed] (A) -- (D);
\draw[dashed] (B) -- (D);
\draw[dashed] (C) -- (D);
\draw[dashed] (A) -- (2.5,1.5) -- (B);

\node[below left] at (O) {$O$};
\node[below] at (A) {$A$};
\node[left] at (B) {$B$};
\node[left] at (C) {$C$};
\node[above right] at (D) {$D$};
\end{tikzpicture}
\caption*{Рисунок 15.2}
\end{figure}

\begin{minipage}{0.5\textwidth}
$\vd = \overline{OA} + \overline{OB} + \overline{OC} =$

$= \alpha \va + \beta \vb + \gamma \vc$

$\Rightarrow \vd$ --- линейно выражается через $\va, \vb, \vc \Rightarrow$

$\Rightarrow \Sigma = \{\va, \vb, \vc, \vd\}$ линейно зависима.
\end{minipage}

\textbf{Пример}

Исследовать на линейную зависимость систему векторов $\Sigma = \{\va,
\vb, \vc\} \subset V_3$, если $\va(2; -1; 3)$; $\vb(3; 2; 1)$; $\vc(1; 0; 4)$.

\textbf{Решение}

$\Sigma = \{\va, \vb, \vc\}$ --- линейно зависима $\Leftrightarrow
\va, \vb, \vc$ --- компланарны $\Leftrightarrow (\va, \vb, \vc) = 0$
$$
(\va, \vb, \vc) =
\begin{vmatrix} 2 & -1 & 3 \\ 3 & 2 & 1 \\ 1 & 0 & 4
\end{vmatrix} =
\begin{vmatrix} 2 & -1 & -5 \\ 3 & 2 & -11 \\ 1 & 0 & 0
\end{vmatrix} =
\begin{vmatrix} -1 & -5 \\ 2 & -11
\end{vmatrix} = 21 \ne 0
$$
$\Longrightarrow$ векторы $\va, \vb, \vc$ линейно независимы.

\end{document}
